\par\noindent\textbf{(1)} 半经典Hamiltonian H=p²/(2m)+m(ω\_x²x²+ω\_y²y²+ω\_z²z²)/2。六维相空间体积∝ε³/(ℏ³ω\_xω\_yω\_z)，求导得激发态密度 g(ε)=ε²/[2ℏ³ω\_xω\_yω\_z]（不计单独基态），其ε²与均匀三维自由气体的√ε不同。
\par\noindent\textbf{(2)} 在临界点化学势趋最低能级0，激发粒子数 N\_ex=∫g(ε)/(e\textasciicircum{}\{βε\}−1)dε=ζ(3)(k\_BT)³/(ℏ³ω\_xω\_yω\_z)。令N=N\_ex(T\_c)，得 k\_BT\_c=ℏ(ω\_xω\_yω\_z)\textasciicircum{}\{1/3\}[N/ζ(3)]\textasciicircum{}\{1/3\}；凝聚相 N₀/N=1−(T/T\_c)³。
\par\noindent\textbf{(3)} 热能 U\_ex=∫εg(ε)/(e\textasciicircum{}\{βε\}−1)dε=3ζ(4)(k\_BT)⁴/(ℏ³ω\_xω\_yω\_z)，基态能可整体平移为零。于是T<T\_c时 C\_V=12ζ(4)k\_B(k\_BT)³/(ℏ³ω\_xω\_yω\_z)，在T\_c下侧有随T³上升的热容。
\par\noindent\textbf{(4)} 谐振陷阱下凝聚分数指数为3，而均匀三维气体指数为3/2；差别源于态密度能量幂次。有限N下能级离散且基态能不为零，修正的相对量级约N\textasciicircum{}\{−1/3\}，若k\_BT并不远大于ℏω\_i则半经典积分失效。
